\section{Limitations and Future Work}
\label{sec:limitations}

\subsection{Limitations}
\begin{description}[leftmargin=1.4em, style=nextline]
  \item[OCR is not perfect.] Recognition errors propagate into search. Because
        native find is exact, a misread character can make a word unfindable by
        its true spelling. Correction variants mitigate the most common
        digit/letter confusions but do not cover the general case.
  \item[Cost scales with the document.] Recognition is CPU-bound and runs on the
        client. Very large or very dense documents take correspondingly longer to
        fully index, even though lazy rendering keeps the document immediately
        readable and caching makes subsequent loads instant.
  \item[Complex layouts can misalign.] Multi-column text, unusual fonts, and
        heavy graphical backgrounds can produce slight overlay drift or
        degraded recognition.
  \item[Device-dependent performance.] Throughput depends on the number of CPU
        cores and the device's overall capability, since the worker pool scales
        with available concurrency.
  \item[Scope is currently PDF-centric.] Image text on ordinary web pages is not
        yet handled; the content script that would address this is presently a
        placeholder.
\end{description}

\subsection{Future Work}
\begin{itemize}
  \item \textbf{Higher OCR accuracy.} Richer post-processing---dictionary-based
        correction, confidence-weighted variants, and language-model
        re-ranking---to reduce unfindable misreads beyond the current
        digit/letter mapping.
  \item \textbf{Better rotation and skew handling.} Automatic deskew before
        recognition for photographed or imperfectly-scanned pages.
  \item \textbf{Broader language support.} Additional bundled language models and
        automatic script detection.
  \item \textbf{Acceleration.} Exploiting GPU/WebAssembly SIMD acceleration and
        preprocessing whole documents in the background to shorten time-to-fully-
        searchable on large files.
  \item \textbf{General web-page image text.} Extending the pipeline, via the
        content script, to make text inside ordinary \code{<img>} elements
        searchable, not just PDFs.
  \item \textbf{Quantitative evaluation.} A benchmark suite measuring recognition
        accuracy (e.g.\ character/word error rate) and per-page latency across a
        corpus of representative scanned documents, to turn the qualitative
        results of Section~\ref{sec:results} into reproducible numbers.
\end{itemize}
